% !TeX program = xelatex
\documentclass[11pt,oneside]{article}
\usepackage{ProblemsetStyle}

\hypersetup{
    pdftitle={20260909 Problem set - Solutions},
    pdfauthor={김정우},
    pdfsubject={Science Quiz mathematics practice problem set solutions},
    pdfcreator={XeLaTeX}
}

\begin{document}

\problemsettitle{20260909 Problem set}

\begin{problemsolution}{1}
배각공식
\[
    (2\cos\theta)^2-2=2\cos(2\theta)
\]
에 의해 귀납적으로
\[
    a_n
    =
    2\cos\left(\frac{2^{n+1}\pi}{7}\right)
\]
을 얻는다. $2$의 거듭제곱을 $14$로 나눈 나머지는 주기 $3$을 가지며,
\[
    2^{2027}\equiv4\pmod{14}.
\]
따라서
\[
    \boxed{a_{2026}=2\cos\left(\frac{4\pi}{7}\right)}.
\]
\end{problemsolution}

\begin{problemsolution}{2}
Vi\`ete's formulas에 의해
\[
    \alpha+\beta+\gamma=0,
    \qquad
    \alpha\beta+\beta\gamma+\gamma\alpha=-3.
\]
또한 $1$은 주어진 다항식의 근이 아니므로 분모는 모두 $0$이 아니다. 통분하면 분자는
\[
    \begin{aligned}
        &(1-\beta)(1-\gamma)
        +(1-\gamma)(1-\alpha)
        +(1-\alpha)(1-\beta)\\
        &\qquad
        =3-2(\alpha+\beta+\gamma)
        +(\alpha\beta+\beta\gamma+\gamma\alpha)
        =0.
    \end{aligned}
\]
따라서
\[
    \boxed{0}.
\]
\end{problemsolution}

\begin{problemsolution}{3}
주어진 식은
\[
    z^2-z+1=0
\]
과 동치이므로
\[
    z=e^{\pm i\pi/3}.
\]
따라서
\[
    z^m+z^{-m}
    =
    2\cos\left(\frac{m\pi}{3}\right).
\]
$2026\equiv4\pmod6$이므로
\[
    \boxed{
        z^{2026}+z^{-2026}
        =
        2\cos\left(\frac{4\pi}{3}\right)
        =-1
    }.
\]
\end{problemsolution}

\end{document}
